\documentclass{article}
\usepackage[a4paper,margin=1in]{geometry}
\usepackage{titlesec}
\titleformat{\section}{\large\bfseries}{\thesection}{1em}{}
\usepackage{booktabs}
\usepackage{tabularx}
\usepackage[utf8]{inputenc}
\usepackage{geometry}
\usepackage{booktabs}
\usepackage{longtable}
\usepackage{amsmath}
\usepackage{amsfonts}
\usepackage{graphicx}
\usepackage{blindtext}
\usepackage{hyperref}
\usepackage{float} % <-- add this to force table positioning
\usepackage{natbib} % <-- NEW: to handle references
\usepackage{setspace}
\usepackage{array}
\usepackage{dcolumn}
\usepackage{threeparttable}
\usepackage{tikz}
\usepackage{amsmath}
\usetikzlibrary{decorations.pathreplacing}
\usepackage{pdflscape} % in your preamble
\usepackage{tabularray}
\usepackage{booktabs}      % for \toprule, \midrule, \bottomrule
\usepackage{multirow}      % for \multirow
\usepackage{array}         % for p{...} column type with better spacing
\usepackage{threeparttable} % for table notes
\usepackage{float}         % for [H] positioning\setcounter{secnumdepth}{2}
\setstretch{1.15}
\setlength{\parskip}{6pt}


\title{Research Proposal: Topics in Econometrics}
\author{Jordi Torres Vallverdú}
\date{\today}

\begin{document}

\maketitle
\vspace{-2em}


\section{Introduction}

In many education systems, students must make binding track choices at a very 
young age. These decisions occur under severe uncertainty about ability, 
preferences, and future success \citep{betts2011economics, dustmann2017returns, 
giustinelli2017evolution}, and they are largely irreversible: early tracking 
placements strongly shape higher education access, labor market trajectories, 
and the dynamics of skill accumulation \citep{hanushek2006does, dustmann2017returns, 
de_groote_dynamic_2025, humphries_joensen_veramendi_2025_complementarities}. 
The information students have when making these choices is therefore consequential 
--- and where that information comes from matters.

A large literature documents the rigidities and inequalities generated by 
tracking, but the ``soft'' institutional mechanisms governing these choices 
remain largely unexplored. In many OECD countries --- including Germany, Belgium, 
the Netherlands, Switzerland, and Austria --- schools issue formal track 
recommendations that act as powerful signals at a stage when students are 
actively forming beliefs about their own comparative advantage 
\citep{piopiunik2015teacher, Papageorge2020Expectations, giustinelli2017evolution}. 
This advice embeds soft information --- effort, relative class rank, classroom 
behavior --- that goes beyond objective test scores, and it is susceptible to 
institutional bias \citep{carlana_implicit_2019, brunello2025teacher}. Whether 
this advice helps students or distorts their choices depends on what schools 
know, what schools want, and how students update beliefs in response.

A natural framework for thinking about this is the growing structural literature 
on educational choice under information frictions 
\citep{stinebrickner2012learning, stinebrickner2014academic, larroucau2024dynamic, 
arcidiacono2025college}. In this literature, students observe noisy signals --- 
grades, wages, output --- and update Bayesianly about their latent ability. 
But signals are treated as features of the environment: they arrive 
\emph{exogenously}. Many real institutional signals are not like this. They 
are issued by strategic actors with their own objectives: teachers, supervisors, 
lenders, doctors. Modeling these signals as exogenous misses what may be the 
central feature of how information flows through educational institutions --- 
that the signal a student receives reflects both what the institution knows 
and what the institution wants.

This project models educational choice under uncertainty when one of the signals 
is endogenously produced by an institution with conflicting interests. Schools 
in our setting play a dual role. They observe more than students do --- relative 
class rank, behavioral cues, accumulated experience mapping past cohorts to 
outcomes --- which makes their recommendations potentially informative. But 
they also face capacity constraints and outcome-reporting incentives that may 
not align with student welfare, which means their recommendations are also 
shaped by institutional pressures unrelated to what is best for the student. 
A school may, in equilibrium, push high-ability students into lower tracks 
to manage cohort sizes. Whether this happens, how it varies across school 
types, and what it costs students are the empirical questions this paper 
addresses. Students, on the other side, do not receive this advice passively: 
they are Bayesians who combine the school's signal with public information 
(grades) and their own private knowledge, with response heterogeneity moderated 
by family background.

Belgium provides a clean setting to study this. At age 14, students commit to 
academic, technical, or vocational tracks; schools issue formal recommendations 
annually; and the Belgian LOSO dataset \citep{de_groote_dynamic_2025} observes 
the full information set schools use to recommend (relative grade rank, relative 
effort rank), the standardized exams that form a balanced grade panel, and a 
dimension of plausibly different school behavior: cycle architecture. Some 
schools offer only the first cycle of secondary education, while others offer 
all three. Schools facing different cohort flows have plausibly different 
incentives in how they recommend --- a source of variation we exploit for the 
counterfactual.

The project aims to make two contributions. First, it extends the structural 
learning literature \citep{stinebrickner2012learning, arcidiacono2025college, 
larroucau2024dynamic, bunting2024heterogeneity} to settings where one of the 
signals is generated by a strategic institutional actor rather than arriving 
exogenously, providing the first such model applied to secondary-school 
tracking. Second, it provides empirical evidence on how 
heterogeneity in school recommendation rules --- driven by institutional 
architecture --- affects student welfare and the SES gradient in track 
outcomes, contributing to literatures on teacher expectations and the role of soft signals in tracking systems \citep{piopiunik2015teacher, carlana_implicit_2019, brunello2025teacher}.
\subsection{Institutional Setting}

Belgium has a secondary education system that implements tracking at a very early age. At 12, students enter high-school and must decide between two general tracks --- denominated \textit{A-stroom} and \textit{B-stroom}. These two tracks, although general, already generate some classification: \textit{A-stroom} prepares for academic and technical tracks, while \textit{B-stroom} leads to vocational tracks.\footnote{The dataset I was given by Olivier De Groote, used in \citet{de_groote_dynamic_2025}, classified students from \textit{A-stroom} and \textit{B-stroom} into the tracks that later will be binding. This was done based on the track --- A leads to TSO and ASO and B to BSO --- and on the electives picked, which were used to distinguish ASO from TSO. For more information see Appendix B of \citet{de_groote_dynamic_2025}.} The decision to enter A-stroom or B-stroom at grade 1 of high-school is generally unconstrained, although it is informed by previous school grades and advice from school teachers.

At age 14, students make a binding track choice. The options are: the academic track (\textbf{ASO}), the technical track (\textbf{TSO}), and the vocational track (\textbf{BSO}). The first is the most common option and leads mostly to university and professional colleges; the second is the least common and leads mostly to professional college; the third is the second most common and leads to professional training. Up to high-school graduation there is little upward mobility (BSO $\to$ TSO $\to$ ASO) but substantial downward mobility. Students may either learn that their skill does not match their track or be forced to downgrade if they fail. At the end of each year, students are issued a general certificate of performance: the A certificate allows students to continue and pick whatever option they want (upgrade, downgrade, or continue), the C certificate forces students to repeat track, and the B certificate forces them to either downgrade or repeat the course in the same track.

Within ASO and TSO, students at age 14 also choose a subtrack tailored to their interests --- e.g., ``Latin and math'' or ``social sciences.'' The subtrack is highly predictive of the field they will enroll in college. Although entry to college is not restricted by entry exams and students are free to choose any field,\footnote{University is mostly public and has very low fees.} students risk skill mismatch and dropout. Accumulated skill in a subtrack during high-school is thus important for a successful higher-education path.

A key institutional feature of the Flemish system is the role schools play in these decisions. At the end of each high-school year, the school issues an official recommendation about which track and subtrack a student should be studying next year. These recommendations are part of the academic evaluation --- which also includes the official certificate --- and are formally communicated to families in June.\footnote{The recommendation is formally issued by the class teacher (\textit{Klasstitularis}) after consultation with colleagues. We refer throughout to ``schools'' as the issuer for generality, but the institutional reality is teacher-based. We treat the recommendation as the equilibrium outcome of within-school deliberation and informal communication with parents during the year; identifying parental influence as a separate structural object is beyond the scope of this paper. CF, however, need to avoid changing information structure to survive Lucas Critique.}

The decision at age 14 is consequential because it can lock students in certain paths. Downward mobility is possible based on certificates and performance, but upward mobility is rare due to skill accumulation; horizontal mobility is possible but movement from non-STEM to STEM is rarely observed. In the following years, schools again re-recommend and issue certificates, but the lock-in effects of the age-14 choice make subsequent recommendations less consequential.

A second key feature is institutional variation across schools. Schools differ in two dimensions: which tracks they offer (ASO, TSO-BSO, or all together --- roughly 33\% each) and which cycles they offer (some only cycle 1, grades 1-2; others cycles 1, 2, and 3). The latter --- cycle architecture --- generates plausibly different incentives for school recommendation behavior, and forms one of the central sources of identifying variation for the counterfactuals I have in mind.

\subsection{Data}

The LOSO dataset follows a cohort of \(\approx\)6,400 students from 89 high-schools in the Flemish region of Belgium, collected during the 1990s. The dataset tracks students from age 12 through higher-education entry --- approximately 8.5 years.\footnote{Schools were selected from two regions: Dendermonde-Zele-Hamme and Diest-Hasselt-Genk-Houthalen-Leopoldsburg-Beringen, based on criteria such as school size, administrative structure, and representation of different school types (public and private education).} Here I provide a summary of the available information, at both the student and the school level:

\paragraph{Student-side variables.}
\begin{itemize}
  \item \textbf{Track choices.} For each year, we observe the track-subtrack chosen, the certificate received, and dropout. Attrition is minimal. The first two years of higher education are also observed, including type of college, degree etc.
  \item \textbf{Grades.} Observed for grades 1, 2, 4, and 6 in both Dutch and Math. Crucially, these grades are administered in a standardized format across the entire sample, independently of the track or subtrack the student is in. This balanced panel of comparable performance measures is what enables the AAMR-style identification of learning parameters without selection correction (\S\ref{sec:model_sketch}).
  \item \textbf{Baseline measures.} Cognitive skills, non-cognitive skills, and parental involvement, observed at $t=0$. Used for identification of unobserved types via Hu--Schennach measurement system (\S\ref{tab:measurement_map}).
  \item \textbf{Subjective beliefs.} For each year corresponding to grades observation, we observe the track-subtrack the student expects to graduate in (i.e "ASO-STEM"). From grade 4 onwards, we also observe the job or higher-education path the student expects. These reported expectations provide direct moments on student beliefs over future track outcomes and can be used in identification.
  \item \textbf{Location.} Student's (home) location, allowing computation of distance to school as exogenous variation.
\end{itemize}

\paragraph{School-side variables.}
\begin{itemize}
  \item \textbf{Type of school.} Tracks offered and cycles offered (cycle architecture). Roughly 33\% of schools fall into each track-offering category. Cycle architecture (C1 vs C123) is the central source of variation for the counterfactual.
  \item \textbf{Capacity constraints.} Number of seats per grade and track. Reported at the beginning of the panel based on the population, not the sampled students.
  \item \textbf{General school statistics:} Also given at the beginning of the panel and for other students, schools also provide the percentage of dropout, the percentage of students going to H.E etc (success measures).
  \item \textbf{Recommendations.} The track-subtrack recommended to each student for different periods in time. We collapse the 20 official ASO subtracks (same for TSO) into STEM/non-STEM categories for tractability; alternative taxonomies are possible, but the literature agrees on checking this dimension \citep{arcidiacono2025college, humphries_joensen_veramendi_2025_complementarities}.
  \item \textbf{School information set.} For each class, we observe how the school rates the student in the within-class performance distribution in Math, Dutch, both for exams and effort given. These rankings form the basis of the school's information set $Z$ in the model (\S\ref{sec:model_sketch}).
  \item \textbf{Advice process variation.} We additionally observe survey responses on how each school structures its advice process --- whether it holds parent meetings during the year, whether it consults students before issuing recommendations, and what information about future tracks it provides. This information is informative about residual school heterogeneity beyond cycle architecture and can be used in robustness analyses.
\end{itemize}

\section{Model sketch}
\label{sec:model_sketch}

I present the sketch of the model. First, I define the timing and the players. Second, I describe the sources of unobserved heterogeneity, the key primitives, and the information structure. Then the production functions and the school side. Estimation details go to the appendix and are mostly sketched for brevity.
 
\subsection{Environment and timing}

\paragraph{Time.} Discrete annual periods $t \in \{0, 1, 2, 3, \ldots, T\}$:
\begin{itemize}
  \item $t=0$: start of cycle 1. Baseline measurements $M_i$ taken; school enrollment $s_i$ realized.
  \item $t \in \{1, 2\}$: cycle 1. End-of-year standardized exams $G_{it}$ administered; school-issued recommendation $R_{it}$ delivered after grades.
  \item $t = 3$: binding initial track choice $d_{i3}$ at start of cycle 2.
  \item $t =4, 5$: cycle 2 and 3. Annual track-refinement decisions $d_{it}$ allowed within a downgrade trajectory. Again $R_{it}$, $G_{it}$ are realized.
  \item $t = T$: terminal outcome $Y^{\mathrm{post}}_i$ (HE entry / labour-market signal etc).\footnote{We could also model higher education entry and STEM degree choice as in \citet{de_groote_dynamic_2025}} is realized.
\end{itemize}

Track movements are restricted: BSO is absorbing\footnote{Justified by the low rate of upgrades from BSO in the data.}, and within-cycle moves follow the downgrade order ASO $\to$ TSO $\to$ BSO. Upgrades are not feasible, but movements on the STEM/non-STEM margin are allowed ---one could also justify movements from nonSTEM to STEM to be blocked (it is also a rare event in the data).

\paragraph{Players.} Students $i = 1, \ldots, N$; schools $s \in \mathcal{S}$, indexed by cycle-architecture type $\tau(s) \in \{\mathrm{C1}, \mathrm{C123}\}$ (cycle-1-only vs.\ full cycle) and type. Track set is defined as:
\begin{equation}
  \mathcal{J} = \{\mathrm{ASO\text{-}S},\ \mathrm{ASO\text{-}N},\ \mathrm{TSO\text{-}S},\ \mathrm{TSO\text{-}N},\ \mathrm{BSO}\}, \qquad |\mathcal{J}| = 5.
  \label{eq:tracks}
\end{equation}



\subsection{Student's problem}

\subsubsection{Sources of heterogeneity}

The model has three sources of unobserved heterogeneity: $\theta^k_i$, a discrete type known to student and school but not to the econometrician; $\theta^u_i$, a track-specific match quality unknown to all agents at $t=0$; and the different idiosyncratic shocks $\varepsilon_{it}$. This setting takes mostly as reference the notation and setting developed in \citet{arcidiacono2025college, bunting2024heterogeneity}.

\paragraph{Known type $\theta^k_i$.} Discrete:
\begin{equation}
  \theta^k_i \in \{1, \ldots, K\}, \qquad \Pr(\theta^k_i = k) = \pi_k, \quad \sum_k \pi_k = 1.
  \label{eq:thetak}
\end{equation}
Each type $k$ encodes a profile over four latent dimensions: cognitive skills, a non-cognitive component, and a family-pressure component, $(\theta^k_{\mathrm{num},k}, \theta^k_{\mathrm{verb},k}, \theta^k_{\mathrm{noncog},k}, \theta^k_{\mathrm{fam},k})$. Identification follows \citet{humphries_joensen_veramendi_2025_complementarities} and \citet{arcidiacono2025college}: I use selection-free baseline measurements with at least three repeated measures per dimension\footnote{Identification additionally requires that the measurement loadings on latent factors are linearly independent — i.e., that no measure is a perfect linear combination of others within a factor.}. Table~\ref{tab:measurement_map} in the appendix maps observed measures to latent factors.

\paragraph{Unknown type $\theta^u_i$.} Track-specific match quality, unknown to all agents at $t=0$:
\begin{equation}
  \theta^u_i = (\theta^u_{i,j})_{j \in \mathcal{J}} \in \mathbb{R}^{5}.
  \label{eq:thetau}
\end{equation}
Following \citet{arcidiacono2025college}, this is the latent component of the outcome the student would experience in track $j$, net of $\theta^k$ and $X$. It is the object students learn about over the panel; it absorbs effort and own realizations of skills within each of the tracks. 

Conditional on $\theta^k$, the prior on $\theta^u$ is Gaussian:
\begin{equation}
  \theta^u_i \mid \theta^k_i = k \ \sim\ \mathcal{N}\big(m_0(k),\ \Sigma_0\big),
  \label{eq:prior_thetau}
\end{equation}
with $\Sigma_0$ positive definite and not restricted to be diagonal (correlated learning across tracks, as in AAMR). The mean $m_0(k)$ is type-specific\footnote{Potentially, it can also be $s$ specific, as students may enter with different confidence in each school.}.

\textbf{This I probably need to change? unsure?}

\paragraph{Observed covariates.} $X_{it}$ are exogenous observables (gender, age, grade level).

\vspace{1em}


Here I outline the two main assumptions on the initial distribution of skills:

\paragraph{Assumption A1 (sorting on observable types).} Initial school assignment $s_i$ may depend on $\theta^k_i$ but $\theta^u_i \perp s_i \mid \theta^k_i$. This is justified because at $t=0$ uncertainty over $\theta^u$ has not been realized.\footnote{Priors over $\theta^u$ may differ across schools (e.g.\ if students hold more positive expectations in some schools). Testable by regressing $\hat m^G_{i,1}$ on school dummies controlling for $\theta^k$. If significant, $m_0$ is made school-specific.}

\paragraph{Assumption A2 (Gaussian conjugacy on grades).} Grade noise is Gaussian and i.i.d.\ across $t$, conditional on $(\theta^k, \theta^u, X)$. This delivers closed-form Gaussian posteriors after grade signals; recommendation updating requires an additional approximation (A8 below). This is a standard assumption in the learning literature \citep{bunting2024heterogeneity}.


\subsubsection{Production function}
The production function is described next.

\textbf{This I need to update with the new specification. I need to add outcome instead of grade, no? cause there are multiple outcomes we may care about.}

For $t \in \{1, 2, 3, 4\}$ and subject $b \in \{\mathrm{Math}, \mathrm{Dutch}\}$ we define:
\begin{equation}
  G_{ibt} \ =\ \mu_b(\theta^k_i, X_{it}) \ +\ \Lambda_b^\top \theta^u_i \ +\ \varepsilon_{ibt},
  \qquad \varepsilon_{ibt} \sim \mathcal{N}(0, \sigma_b^2),\ \mathrm{i.i.d.}
  \label{eq:grade}
\end{equation}
\begin{itemize}
  \item $\mu_b(\theta^k_i, X_{it})$: type-specific intercept plus linear in $X_{it}$.
  \item $\Lambda_b \in \mathbb{R}^{5}$: loading vector. Math grades load on STEM components, Dutch on non-STEM. Normalize one entry of $\Lambda_b$ per subject for scale.
\end{itemize}

\textbf{On grades, it is true this remains:}

A useful difference from \citet{arcidiacono2025college} is that $G_{ibt}$ are standardized exams observed for every student at every $t$, regardless of school or future track. Their grade likelihood requires participation indicators because they only observe grades for the major chosen. Our case does not:
\begin{equation}
  L^G_i \ =\ \prod_{t=1}^4 \prod_{b \in \{\mathrm{Math}, \mathrm{Dutch}\}} \phi\!\big(G_{ibt};\, \mu_b(\theta^k_i, X_{it}) + \Lambda_b^\top \theta^u_i,\, \sigma_b^2\big).
  \label{eq:grade_lik}
\end{equation}
The balanced grade panel sharply identifies $(\Lambda_b, \sigma_b^2, \Sigma_0)$ via cross-individual covariances, without the joint outcome-and-choice likelihood \citet{arcidiacono2025college} require. 


\subsection{Student's problem}

\subsubsection{School choice}
The students enter High school and face subset X if schools, They choose school. Then initial ability is conditional on that. We 

static reduced form?

\subsection{Track choice}

This is the dynamic





\subsection{School's information set and the DGP for $Z$}

One of the key assumptions of the paper is that the school observes more than the student In addition to $(\theta^k_i, G_{it}, X_{it})$, it observes
\begin{equation}
  Z_{it} \ =\ (\mathrm{rank}^{Gb}_{it},\ \mathrm{rank}^{Eb}_{it})_{b \in \{\mathrm{Math}, \mathrm{Dutch}\}},
  \label{eq:Z_observed}
\end{equation}
where $\mathrm{rank}^{Gb}_{it}$ is class-relative rank in grades for subject $b$ and $\mathrm{rank}^{Eb}_{it}$ is class-relative rank in effort for subject $b$. Importantly, we also observe these variables that schools use to make decisions. 

I model these rankings as discretized realizations of a latent school-side signal:
\begin{equation}
  Z^*_{ibt} \ =\ \mu_b(\theta^k_i, X_{it}) + \lambda_b^\top \theta^u_i + \xi_{ibt}, \qquad \xi_{ibt} \sim \mathcal{N}(0,1),
  \label{eq:Zstar}
\end{equation}
with the observed category generated by thresholds:
\begin{equation}
  Z_{ibt} = m \quad \Longleftrightarrow \quad \kappa_{m-1,b} < Z^*_{ibt} \le \kappa_{m,b},
  \label{eq:thresholds}
\end{equation}
for $-\infty = \kappa_{0,b} < \kappa_{1,b} < \dots < \kappa_{M_b,b} = +\infty$. This is standard ordered probit and is useful because it simplifies a bit the procedure we will use to turn the discrete signals into normality. 

Importantly, $Z^*_{ibt}$ has the same structural form as grades \eqref{eq:grade} but with independent noise $\xi_{ibt}$, providing an additional draw on the same underlying $\theta^u$. The recommendation is informative beyond grades because $\xi$ is independent of $\varepsilon$. We can be more precise with the following assumption:

\paragraph{Assumption A4 (signal exhaustion).} $\xi$ is the only school-side draw on $\theta^u$. Equivalently, conditional on $(\theta^k, G, Z)$, no other school-side variable (informal parent--teacher communication, sibling outcomes\dots) is correlated with $\theta^u$. This is the substantive content of treating $Z$ as the school's information set; without it, the recommendation's informational content would be biased.\footnote{This is the main limitation of the model so far. Recommendations partly emerge from informal meetings with parents that I do not observe directly, although $\theta^k_{\mathrm{fam}}$ absorbs parental pressure as a baseline type. The implication is that counterfactuals that change the \emph{information structure} are vulnerable to Lucas critique: if in reality schools and parents keep communicating, a CF that changes these primitives says little about reality. The CFs in \S\ref{sec:cf} are chosen to change incentives or what students see, not the school--family communication channel.}

\subsection{Equilibrium recommendation rule}

In this first version, I do not model the school's strategy structurally\footnote{This part is still work in progress, it requires another layer of work}. The recommendation $R_{it}$ is the equilibrium outcome of school behavior, conditional on the school's information set and capacity constraints. I treat the equilibrium rule $\rho_\tau$ as a primitive of the institutional environment.

Define $R_{it} \in \mathcal{R} = \mathcal{J} \cup \{\emptyset\}$, $|\mathcal{R}| = 6$. Schools commit to a recommendation rule:
\begin{equation}
  \rho_{\tau} \ :\ (\theta^k, G_t, Z_t, X_t, K_t) \ \longmapsto\ \Delta(\mathcal{R}),
  \label{eq:rho_def}
\end{equation}
indexed by school cycle-architecture type $\tau$. Capacity $K_{s,j,t}$ enters $W_{it}$ as an observable covariate that shifts recommendations under tighter constraints. We can parametrise this as multinomial logit:
\begin{equation}
  \Pr\!\big(R_{it} = r \,\big|\, W_{it};\ \tau,t\big)
  \ =\ \frac{\exp(W_{it}^\top \beta_{r,\tau,t})}{\sum_{r' \in \mathcal{R}} \exp(W_{it}^\top \beta_{r',\tau,t})},
  \label{eq:rho_mnl}
\end{equation}
with $W_{it} = (\theta^k_i, G_{it}, Z_{it}, X_{it}, K_{s,j,t})$.

I impose these two assumptions on the school side:

\paragraph{Assumption A5 (rule exclusion).}
\begin{equation}
  R_{it} \ \perp\ \theta^u_i \ \big|\ (\theta^k_i, G_{it}, Z_{it}, K_{s,j,t}, \tau, t).
  \label{eq:A5}
\end{equation}
$\theta^u$ is unobserved by everyone, so the rule cannot condition on it directly. The rule conditions on what the school observes, and the recommendation's informativeness about $\theta^u$ flows entirely through $p(Z \mid \theta^u, \cdot)$. Together with A4, this pins down $R$ as a deterministic function of the school's information set up to the logit shock.

\paragraph{Assumption A6 (cycle-architecture).} The equilibrium rule depends on schools only through $\tau(s)$. Testable by comparing within-$\tau$ to between-$\tau$ variation in $\hat\rho_s$. If within-$\tau$ heterogeneity is large, the CF in \S\ref{sec:cf} is reinterpreted as replacing C123 rules with the C1-average rule.


\subsection{Alternative school rule}
The school issues a recommendation $R_{it} \in \mathcal{R}$ to each student at the end of each period $t$. I decompose this into two stages, motivated by the data: capacity constraints bind at the track level (ASO/TSO/BSO), where I observe seats per track, but not at the subtrack level (STEM/nonSTEM), where student preferences appear to be the binding margin.

\subsubsection{Stage 1: Track recommendation (capacity-constrained)}

For each student $i$ and track $j \in \{ASO, TSO, BSO\}$, the school constructs a track-specific academic index:
\begin{equation}
  I_{ij}^{\mathrm{track}} = W_{it}^\top \alpha_{j,\tau} + \nu_{ijt}, \qquad \nu_{ijt} \sim \text{T1EV},
  \label{eq:track_index}
\end{equation}
where $W_{it} = (\theta^k_i, G_{it}, Z_{it}, X_{it})$ is the school's full information set and $\alpha_{j,\tau}$ are track- and cycle-architecture-specific index weights. The school ranks all students on these indices and recommends the top $x_j^s$ fraction to each track, where
\begin{equation}
  x_j^s = (1 + b_s) \cdot \frac{K_{jt}^s}{N_s},
  \label{eq:capacity_share}
\end{equation}
$K_{jt}^s$ is the observed capacity for track $j$ at school $s$, and $b_s$ is a school-specific compliance buffer reflecting that not all students follow the recommendation or that capacity can be expanded in a given school exceptionally.

The track recommendation is:
\begin{equation}
  R_i^{\mathrm{track}} = \arg\max_{j \in \{ASO, TSO, BSO\}} \left\{  I_{ij}^{\mathrm{track}} - c_{jst} \right\},
  \label{eq:track_rec}
\end{equation}
where $c_{jst}$ is an endogenous cutoff pinned down by the capacity-clearing condition: the expected share of students recommended to track $j$ must equal $x_j^s$. \textbf{Doubt, how to formalize this... how to identify the cutoff exactly. How to estimate. I think this captures the idea of the ranking... because it is all relative.}


\textcolor{blue}{\textbf{//! to complete}} also pending to check A/B/C certificate how enters the decision first. Need to check if A/B/C are captured by capacity directly.

\subsubsection{Stage 2: Subtrack recommendation (unconstained)}

Conditional on $R_i^{\mathrm{track}} = j$ for $j \in \{ASO, TSO\}$, the school assesses comparative advantage along the STEM/nonSTEM dimension:
\begin{equation}
  R_i^{\mathrm{sub}} = \text{STEM} \iff W_{it}^\top \gamma_{\text{S},\tau} + \omega_{it} > W_{it}^\top \gamma_{\text{N},\tau} + \omega'_{it},
  \label{eq:sub_rec}
\end{equation}
where $\omega, \omega'$ are T1EV. No capacity constraint enters this stage. The key covariates are the $Z$ rankings: relative math rank drives STEM recommendations, relative Dutch rank drives nonSTEM, consistent with the empirical finding that schools recommend on comparative advantage at the subtrack level.

\subsubsection{The school's objective (Stackelberg)}
\textcolor{blue}{\textbf{//! to complete, still shacky}}

Potentially, the index weights $(\alpha_{j,\tau}, \gamma_\tau)$ are not arbitrary --- they are the school's optimal choice, anticipating how students will respond. The school solves a per-period problem:
\begin{equation}
  \max_{\alpha_\tau, \gamma_\tau} \sum_{i=1}^{N_s} \sum_{j \in \mathcal{J}} P\!\big(d_{it} = j \,\big|\, R_{it}(\alpha, \gamma),\, s_{it}\big) \cdot \pi_{ij\tau},
  \label{eq:school_obj}
\end{equation}
subject to:
\begin{equation}
  \sum_{i} P\!\big(R_i^{\mathrm{track}} = j \,\big|\, \alpha_\tau\big) = x_j^s \qquad \forall\, j \in \{ASO, TSO, BSO\}.
  \label{eq:cap_constraint}
\end{equation}

The student's choice probability $P(d_{it} = j \mid R_{it}, s_{it})$ comes from the estimated dynamic model (\S\ref{sec:beliefs}--\ref{eq:ccp}). The school takes this response function as given\footnote{The issue in CF is if we change the policy, this changes CCPs etc... I don't know how hard the whole thing becomes.}. Changing $\alpha_\tau$ changes the recommendation, which changes the student's posterior beliefs (through the moment-matching update in \S\ref{sec:beliefs}), which changes the student's track choice, the grades etc.

The school's per-student payoff adapts \citet{fu_mehta}:
\begin{equation}
  \pi_{ij\tau} = \mathbb{E}^{school}\!\big[G_{ij,t+1} \,\big|\, \theta^k_i, G_{it}, Z_{it}\big] + \omega_\tau \cdot P\!\big(G_{ij,t+1} > y^* \,\big|\, \theta^k_i, G_{it}, Z_{it}\big),
  \label{eq:school_payoff}
\end{equation}
where the expectation is under the school's posterior (which conditions on $Z$, unlike the student's). Here schools may only care about next period results... but ofc that depends on the expected ability per track they think students have... this needs to be more precise


\subsubsection{Equilibrium}
\textcolor{blue}{\textbf{//! to complete}}

\subsection{Student's beliefs and Bayesian updating}
\label{sec:beliefs}

Let $\mathcal{H}_{it}$ be the student's information at end of period $t$. Within each $t \in \{1,\ldots,4\}$, updating happens in two stages. Here I use the fact that students first observe grades of the standardized exams and that recommendations only come at the end of the year:

\paragraph{Stage 1: posterior after grades.} Given $\mathcal{H}_{i,t-1} \sim \mathcal{N}(m^H_{i,t-1}, \Sigma^H_{i,t-1})$ and observed grades $G_{it}$, conjugacy gives
\begin{equation}
  \theta^u_i \,\big|\, \mathcal{H}_{i,t-1}, G_{it} \ \sim\ \mathcal{N}(m^G_{it}, \Sigma^G_{it}),
  \label{eq:post_grade}
\end{equation}
with
\begin{align}
  \big(\Sigma^G_{it}\big)^{-1} &\ =\ \big(\Sigma^H_{i,t-1}\big)^{-1} + \sum_b \frac{\Lambda_b \Lambda_b^\top}{\sigma_b^2}, \label{eq:Sigma_grade}\\
  m^G_{it} &\ =\ \Sigma^G_{it}\!\left[\big(\Sigma^H_{i,t-1}\big)^{-1} m^H_{i,t-1} + \sum_b \frac{\Lambda_b\big(G_{ibt} - \mu_b(\theta^k_i, X_{it})\big)}{\sigma_b^2}\right]. \label{eq:m_grade}
\end{align}

This is a standard result in the literature \citet{arcidiacono2025college, bunting2024heterogeneity}. Now, the second stage is slightly more tricky ---given the discrete nature of the recommendations---, I propose the following steps, although this is still preliminary: 


\paragraph{Stage 2: posterior after recommendation.} Given the post-grade 
posterior and observed $R_{it}$, Bayes gives
\begin{equation}
  f(\theta^u \mid \mathcal{H}_{i,t-1}, G_{it}, R_{it}, \tau)
  \propto \phi(\theta^u; m^G_{it}, \Sigma^G_{it}) \cdot L_i(R_{it} \mid \theta^u; \tau),
  \label{eq:post_rec}
\end{equation}
where the perceived recommendation likelihood, derived under A5 by 
integrating out $Z$, is
\begin{equation}
  L_i(R \mid \theta^u; \tau)
  = \sum_{m=1}^{M} \rho_\tau(R \mid W_{-Z}, Z = m) \cdot 
    \big[\Phi(\kappa_m - \mu - \lambda^\top \theta^u) - 
         \Phi(\kappa_{m-1} - \mu - \lambda^\top \theta^u)\big].
  \label{eq:rec_likelihood}
\end{equation}
The integral over the latent $Z^*$ has closed form via $\Phi$; the sum is 
finite over the ordered-probit cells.

\paragraph{Assumption A8 (Gaussian moment-matching).} The post-recommendation
posterior is approximated as Gaussian by matching the first two moments of
the exact mixture-of-truncated-Gaussians posterior. This is the assumed-density-filtering step of \citet{minka2001expectation}; \citet{sarkka2013bayesian} gives the textbook treatment in the filtering literature\footnote{Work in progress... approximation error should be low, but need to go step by step. This should be feasible and computationally very simple.}:
\begin{equation}
  \theta^u_i \mid \mathcal{H}_{i,t-1}, G_{it}, R_{it} 
  \stackrel{\mathrm{approx}}{\sim} \mathcal{N}(m^R_{it}, \Sigma^R_{it}),
\end{equation}
with closed-form expressions
\begin{align}
  m^R_{it} &= \sum_{m} \pi^R_m \cdot \mathbb{E}_m[\theta^u], \\
  \Sigma^R_{it} &= \sum_m \pi^R_m \big[V_m + (\mathbb{E}_m[\theta^u] - m^R_{it})
                                           (\mathbb{E}_m[\theta^u] - m^R_{it})^\top\big],
\end{align}
where $\pi^R_m \propto \rho_\tau(R \mid Z=m) \cdot \Pr_m$ are the recommendation-
weighted cell probabilities, $\mathbb{E}_m[\theta^u]$ and $V_m$ are the conditional 
mean and variance of $\theta^u$ given $\{Z = m\}$ (closed form via inverse-Mills 
ratios for doubly-truncated Gaussians), and $\Pr_m$ are the marginal cell 
probabilities under the post-grade prior. Full derivations in Appendix A.

\paragraph{Intuition behind this result:} The student doesn't see $Z$, but they see $R$, the school's discrete summary of it. When a student receives ASO-STEM, they reason 
backward: the school must have seen $Z$ values consistent with that recommendation, 
given the student's grades. That backward inference is what shifts beliefs.

How much beliefs move depends on the school's rule. If a school recommends 
ASO-STEM to nearly everyone with similar grades, the recommendation adds little 
beyond grades and beliefs barely move. If the school reserves ASO-STEM for 
specific $Z$ patterns (high effort rank, top of class), receiving ASO-STEM is a 
strong signal and beliefs move sharply. Within tracks, beliefs shift in the 
direction the recommendation suggests: a high recommendation pulls the posterior 
mean of $\theta^u$ upward for the recommended track\footnote{The data supports that students reach on recommendations even conditional on standardized grades results. This supports this approach.}.

This is the channel through which institutional incentives translate into 
student outcomes. A school facing tight capacity may issue ASO-STEM less 
frequently than a school with the same students but slack capacity. ``ASO-STEM'' 
then carries different information at the two schools --- not because students 
learn differently, but because the schools \emph{say} different things in 
equivalent situations. Students Bayesianly absorb whatever the rule produces; 
the counterfactuals in \S\ref{sec:cf} vary the rule.

\subsection{State, action, and dynamic choice}

The state of student $i$ at the start of $t$ is
\begin{equation}
  s_{it} \ =\ \big(\theta^k_i,\ m^H_{i,t-1},\ \Sigma^H_{i,t-1},\ d_{i,t-1},\ X_{it}\big).
  \label{eq:state}
\end{equation}
$\Sigma^H_{i,t-1}$ is deterministic given the past sequence of actions (Kalman variance updates do not depend on realized signals), so the effective state is $(\theta^k_i, m^H_{i,t-1}, d_{i,t-1}, X_{it})$. $d_{i,t-1}$ captures the previous track and switching costs.

The action set at $t$ is $\mathcal{A}_t$. At $t=3$, $\mathcal{A}_3 = \mathcal{J}_i \subseteq \mathcal{J}$ restricted by the certificate (A/B/C). For $t \geq 4$, $\mathcal{A}_t$ is the set of feasible downgrades from $d_{i,t-1}$ (BSO absorbing).

Flow utility for $a \in \mathcal{A}_t$:
\begin{equation}
  u(s_{it}, a; \theta) \ =\ \alpha_a + \zeta_a^\top \theta^k_i + \delta_{\mathrm{SES}(i)} \cdot e_a^\top m^H_{i,t-1} + \gamma^\top X_{it} + c(d_{i,t-1}, a) + \eta_{ita},
  \label{eq:flow}
\end{equation}
where $\zeta_a$ are track-specific loadings on the discrete type, $e_a$ is the $a$-th unit vector (so $e_a^\top m^H_{i,t-1}$ is the posterior mean of $\theta^u_{i,a}$), $c(\cdot, \cdot)$ captures switching costs, $\eta_{ita} \sim$ T1EV i.i.d., and $\delta_{\mathrm{SES}(i)}$ is an SES-specific weight on belief means\footnote{Key heterogeneity in our mdoel would be parental SES and priors. Here baseline beliefs can be made dependent of family priors + how students update can also potentially be a function of SES. Idea is that what matters in all these cases of informational asymmetries is the background: high-SES students likely have more information at home or stronger prior; while low-SES students may be more impacted by the policy.}. Integration over $\theta^k$ uses EM weights $q_{i,k}$.\footnote{Note that in this first proposal I have only modeled the effect of recommendation on beliefs. However, there are other channels through which recommendations can affect actions. One is authority; it does not change beliefs but students just follow what schools decide. Including an effect of advice directly on the flow utility ---together with the effect on belief--- would isolate this. I leave this also for future work.}

The Bellman:
\begin{equation}
  V_t(s_{it}) \ =\ \mathbb{E}_\eta \max_{a \in \mathcal{A}_t} \Big\{ u(s_{it}, a; \theta) + \beta\, \mathbb{E}\big[V_{t+1}(s_{i,t+1}) \,\big|\, s_{it}, a\big] \Big\},
  \label{eq:bellman}
\end{equation}
with T1EV log-sum:
\begin{equation}
  \bar V_t(s_{it}) \ =\ \log \sum_{a \in \mathcal{A}_t} \exp(v_t(s_{it}, a)), \qquad v_t(s_{it}, a) \ =\ u(s_{it}, a; \theta) + \beta\, \mathbb{E}\big[\bar V_{t+1}(s_{i,t+1}) \,\big|\, s_{it}, a\big].
  \label{eq:logsum}
\end{equation}
Choice probabilities:
\begin{equation}
  \Pr(d_{it} = a \mid s_{it}) \ =\ \frac{\exp v_t(s_{it}, a)}{\sum_{a' \in \mathcal{A}_t} \exp v_t(s_{it}, a')}.
  \label{eq:ccp}
\end{equation}

\textcolor{blue}{\textbf{//! to complete, revise inclusion of delta}}

\paragraph{Assumption A11 (finite dependence).} The downgrade-only dynamics with absorbing BSO yield finite dependence\footnote{Or dropout, more precisely}: from any current track $a$, a path of length $K = T - t$ ending in BSO has identical conditional state distribution as the reference path "drop to BSO immediately" after $K$ periods, \citet{arcidiacono_miller_2011, arcidiacono2025college} inversion then expresses differenced FV terms in terms of CCPs on the post-recommendation belief state, computed flexibly in Stage 3 of estimation. This avoids backward recursion and full solution methods on the dynamic model. 


\subsection{Counterfactuals}
\label{sec:cf}

The model has two economic forces that pull in opposite directions. On the one hand, the school knows more than the student: $Z$ contains residual information about $\theta^u$ that students cannot access on their own, so the recommendation is informative. On the other hand, the rule that maps $Z$ into $R$ is shaped by capacity and cycle architecture, so it is also \emph{distorted}. The recommendation is at once a useful signal and a sorting device. Each counterfactual below isolates one side of this tension. I keep the institutional architecture --- who talks to whom, how recommendations are produced --- fixed, in line with A4 and the partial-equilibrium scope. Some ideas for CF in this setting are: 

\paragraph{CF1: information disclosure.} Show students $Z$ directly while leaving $\rho_\tau$ unchanged. Students update on the school's underlying signal as if it were just another grade; the recommendation $R$ becomes redundant. This isolates the \emph{value of the recommendation as a signal} from its role as a sorting device. The non-obvious result is that disclosure need not help: if the school's rule is well-calibrated, $R$ is already a sufficient summary of $Z$ for tracking purposes, and disclosure mostly adds noise to high-SES students who already have good priors. The students who benefit are those whose families update most aggressively on new information --- typically lower-SES, by the $\delta_{\mathrm{SES}}$ channel.

\paragraph{CF2: capacity-only.} Hold $\hat\rho_\tau$ and the rule coefficients $\beta$ fixed; raise C123 capacity to C1 levels ---or make schools C123 recommend like C1, this would be cleaner---. The CF measures how much of the welfare gap between school types is attributable to the binding capacity constraint, separately from any difference in how the two types of schools \emph{would} advise under slack capacity. The non-obvious result is that capacity may matter most for the \emph{marginal} student --- those whom C123 schools nudge into TSO/BSO to manage cohort sizes --- and the welfare effect runs through dynamic skill accumulation rather than through the static track choice.

Ideally, the data also supports modelling the school side structurally. The 
most policy-relevant counterfactuals --- changing the school's objective 
directly, i.e\ aligning incentives across cycle architectures or imposing 
truthful disclosure --- require identifying the school's objective $\omega_\tau$, likely from variation in capacity or institutional rules. \citet{fu_mehta} provides 
a natural template, but I leave this for future iterations.

\section*{Appendix}


\subsection{Measurements avaialble at baseline}
In the following table I map the measures at baseline that I have and how these can be used to recover latent traits:

\begin{table}[H]
\centering
\caption{Mapping of observed measures to latent factors}
\label{tab:measurement_map}
\begin{threeparttable}
\begin{tabular}{p{3.5cm} p{6cm} p{2cm} p{2cm}}
\toprule
\textbf{Latent factor} & \textbf{Observed measures} & \textbf{Source}  \\
\midrule

\multirow{3}{=}{$\theta^k_{\mathrm{num}}$ \\ Numeric ability} 
  & GETLOV numeric reasoning subscore & Student \\
  & Results Math exam School   & Student \\
  & Spatial IQ test  & Student \\
\midrule

\multirow{3}{=}{$\theta^k_{\mathrm{verb}}$ \\ Verbal ability} 
  & GETLOV verbal reasoning subscore & Student \\ 
  & Results Dutch exam School & Student \\
  & Spatial IQ test & Student \\
\midrule

\multirow{4}{=}{$\theta^k_{\mathrm{noncog}}$ \\ Non-cognitive skills} 
  & PMTK achievement motivation & Student  \\
  & PMTK task orientation       & Student  \\
  & PMTK positive test anxiety  & Student \\
  & Big Five subscales (Eindvragenlijst, year 1) & Student \\
\midrule

\multirow{4}{=}{$\theta^k_{\mathrm{fam}}$ \\ Family pressure} 
  & Mother's education          & Parent  \\
  & Father's education          & Parent  \\
  & Parental homework help      & Student  \\
  & Parental income             & Parent  \\
  & Parental influence in tracking decisions             & Student \\

\bottomrule
\end{tabular}
\begin{tablenotes}
\footnotesize
\item \textit{Notes:} Each latent factor is identified from at least three measurements following \citet{hu2008instrumental} and \citet{humphries_joensen_veramendi_2025_complementarities, arcidiacono2025college}. Baseline measurements ($t=0$) are taken before secondary-school enrollment.
\end{tablenotes}
\end{threeparttable}
\end{table}





\subsection{Estimation}

Here I sketch a bit how the estimation would go. I think this is just adapting \citet{arcidiacono2025college} with some simplifications and with the added complication of the nonGaussian signal.

\paragraph{Stage 1 — measurement system + auxiliary regressions $\to q_{i,k}$.} Flexible parametric system: (i) baseline measurements $M_i$ for $\theta^k$ identification (Hu--Schennach); (ii) panel auxiliary regressions on grades, recommendations, choices, and the terminal outcome with type dummies. Output: $q_{i,k}$ and $\pi_k$.

\paragraph{Stage 2 — structural learning parameters.} Weighted EM on the grade equation only, as above. Output: $(\hat\Lambda_b, \hat\sigma_b^2, \hat\mu_b, \hat\Sigma_0)$ and $(m^G_{it}, \Sigma^G_{it})$.

\paragraph{Stage 2.5 — DGP of $Z$ and recommendation rule.} (a) Estimate ordered-probit parameters $(\lambda_b, \kappa_b)$ for $Z$ given Stage-2 posterior moments. (b) Estimate $\beta_{r,\tau,t}$ in \eqref{eq:rho_mnl} by weighted MLE.

\paragraph{Stage 3 — flexible CCPs and FV terms.} Flexible MNL CCPs on the post-recommendation belief state. Differenced FV terms via finite-dependence path to BSO.

\paragraph{Stage 4 — structural flow utility.} Weighted MNL on \eqref{eq:flow} with CCP-derived offset terms. Recovers $(\zeta_a, \delta_{\mathrm{SES}}, \gamma, c, \alpha_a)$.

\paragraph{Identification of CCPs conditional on type.} Given consistent $q_{i,k}$ from Stage 1, type-conditional CCPs are weighted averages \citep{arcidiacono_miller_2011}:
\begin{equation*}
  \widehat{\Pr}\!\big(d_{it} = a \mid s_{it} \approx s, \theta^k = k\big)
  \ =\ \frac{\sum_i q_{i,k}\,\mathbb{1}[d_{it}=a, s_{it} \approx s]}{\sum_i q_{i,k}\,\mathbb{1}[s_{it} \approx s]}.
\end{equation*}



\subsection{Moment-matched post-recommendation posterior}
\label{app:moment_match}

This appendix states the closed-form expressions for $(m^R_{it}, \Sigma^R_{it})$ 
in equation~\eqref{eq:post_rec}. The construction follows the assumed-density-filtering 
approach of \citet{minka2001expectation} adapted to the ordered-probit signal 
structure. I drop $i, t$ subscripts.

\paragraph{Setup.} The post-grade posterior is $\theta^u \sim \mathcal{N}(m^G, \Sigma^G)$. 
The latent school signal $Z^* = \mu + \lambda^\top \theta^u + \xi$ with 
$\xi \sim \mathcal{N}(0,1)$ is observed in $M$ ordered cells via thresholds 
$\kappa_0 < \kappa_1 < \cdots < \kappa_M$. The exact post-recommendation posterior 
is a finite mixture of truncated Gaussians, one per $Z$-cell:
\begin{equation}
  f(\theta^u \mid R) \ \propto\ \phi(\theta^u; m^G, \Sigma^G) \cdot 
  \sum_{m=1}^{M} \rho_\tau(R \mid Z=m) \cdot 
  \big[\Phi(\kappa_m - \mu - \lambda^\top \theta^u) - \Phi(\kappa_{m-1} - \mu - \lambda^\top \theta^u)\big].
  \label{eq:exact_post}
\end{equation}

\paragraph{Matched moments.} Let $\mathbb{E}_m[\theta^u]$ and $V_m$ denote the 
conditional mean and variance of $\theta^u$ within cell $m$ (closed-form in 
$\Phi$, $\phi$ via standard truncated-Gaussian formulas). Let $\Pr_m$ be the 
marginal cell probability under the prior, and define the recommendation-reweighted 
weights
\begin{equation}
  \pi^R_m \ =\ \frac{\rho_\tau(R \mid Z=m) \cdot \Pr_m}{\sum_{m'} \rho_\tau(R \mid Z=m') \cdot \Pr_{m'}}.
\end{equation}
The first two moments of \eqref{eq:exact_post} are
\begin{align}
  m^R &= \sum_{m=1}^M \pi^R_m \cdot \mathbb{E}_m[\theta^u], \label{eq:mR}\\
  \Sigma^R &= \sum_m \pi^R_m \cdot V_m + \sum_m \pi^R_m \cdot (\mathbb{E}_m[\theta^u] - m^R)(\mathbb{E}_m[\theta^u] - m^R)^\top. \label{eq:SR}
\end{align}
The Gaussian moment-matching step \citep{minka2001expectation, sarkka2013bayesian} 
approximates \eqref{eq:exact_post} by $\mathcal{N}(m^R, \Sigma^R)$, exact in the 
first two moments. Validation against numerical quadrature can also be done. 

\bibliographystyle{apalike}
\bibliography{references}

\end{document}

